\documentclass[11pt]{article}
\usepackage[a4paper,margin=0.82in]{geometry}
\usepackage{amsmath,amssymb,graphicx,booktabs,tabularx,array,hyperref,xcolor,enumitem,float,titlesec}
\IfFileExists{microtype.sty}{\usepackage{microtype}}{}
\definecolor{navy}{HTML}{17365D}
\hypersetup{colorlinks=true,linkcolor=navy,citecolor=navy,urlcolor=navy,
pdftitle={Fundamental Signals and Portfolio Construction in US Equities},pdfauthor={Team 34}}
\titleformat{\section}{\Large\bfseries\color{navy}}{\thesection}{0.6em}{}
\titleformat{\subsection}{\large\bfseries}{\thesubsection}{0.5em}{}
\setlist[itemize]{leftmargin=1.35em,itemsep=2pt,topsep=3pt}
\setlength{\parindent}{0pt}\setlength{\parskip}{5pt}
\setlength{\headheight}{16pt}
\makeatletter
\def\ps@teamreport{%
  \def\@oddhead{\vbox{\hbox to\textwidth{\small DSA5205 -- Project 2\hfil Team 34}\vskip4pt\hrule height0.4pt}}%
  \let\@evenhead\@oddhead
  \def\@oddfoot{\hfil\thepage\hfil}\let\@evenfoot\@oddfoot}
\long\def\@makecaption#1#2{%
  \vskip\abovecaptionskip
  \sbox\@tempboxa{\small\textbf{#1:} #2}%
  \ifdim\wd\@tempboxa>\hsize
    {\small\textbf{#1:} #2\par}%
  \else
    \global\@minipagefalse\hb@xt@\hsize{\hfil\box\@tempboxa\hfil}%
  \fi\vskip\belowcaptionskip}
\makeatother
\pagestyle{teamreport}
\newcommand{\nbfig}[3]{\begin{figure}[H]\centering\includegraphics[width=#2\textwidth]{figures/#1.pdf}\caption{#3}\label{fig:#1}\end{figure}}
\newsavebox{\tablebox}
\newcommand{\nbtable}[3]{\begin{table}[H]\centering\caption{#2}\label{tab:#1}\small\sbox{\tablebox}{\input{tables/#1.tex}}\ifdim\wd\tablebox>\linewidth\resizebox{\linewidth}{!}{\usebox{\tablebox}}\else\usebox{\tablebox}\fi\par\vspace{3pt}\begin{minipage}{\textwidth}\footnotesize #3\end{minipage}\end{table}}
\setcounter{tocdepth}{1}
\newcommand{\SelectedName}{Six-factor score / Capitalization}
\newcommand{\SelectedCase}{fundamental\_cap}
\newcommand{\SelectedCAGR}{24.92}
\newcommand{\SelectedVol}{16.81}
\newcommand{\SelectedSharpe}{1.407}
\newcommand{\SelectedMDD}{-20.76}
\newcommand{\SelectedTurnover}{0.92}
\newcommand{\SPYCAGR}{22.19}
\newcommand{\ActiveCAGR}{2.73}
\newcommand{\FitCount}{396}
\newcommand{\CodeCells}{13}
\newcommand{\DataCodeCells}{10}
\newcommand{\LedgerError}{7.54e-10}
\newcommand{\WQList}{WQ004, WQ052, WQ019, WQ035, WQ038, WQ026}
\newcommand{\ResidualAlpha}{7.92}
\newcommand{\ResidualLow}{2.13}
\newcommand{\ResidualHigh}{13.71}
\newcommand{\ResidualT}{2.681}
\newcommand{\RidgeEvalCAGR}{12.31}
\newcommand{\TreeEvalCAGR}{15.72}
\newcommand{\SelectedCILow}{-8.79}
\newcommand{\SelectedCIHigh}{16.40}
\newcommand{\CostTwentyCAGR}{24.34}
\newcommand{\CostThirtyCAGR}{23.77}
\newcommand{\WithoutFourCAGR}{25.33}
\newcommand{\WithoutFourSharpe}{1.271}
\newcommand{\SelectionLink}{3.72}
\newcommand{\WeightingLink}{9.55}
\newcommand{\UniverseLink}{-10.54}
\newcommand{\FirstYearCAGR}{34.96}
\newcommand{\FirstYearSPY}{26.62}
\newcommand{\SecondYearCAGR}{15.56}
\newcommand{\SecondYearSPY}{17.87}
\newcommand{\MarketLinked}{48.76}
\newcommand{\StyleLinked}{-13.00}
\newcommand{\SpecificLinked}{20.80}
\newcommand{\CostLinked}{-1.16}

\begin{document}
\thispagestyle{empty}
\begin{center}
{\Large National University of Singapore\par}
\vspace{0.2cm}{\large DSA5205 -- Project 2\par}
\vspace{0.6cm}
{\LARGE\bfseries\color{navy} Fundamental Signals and Portfolio Construction in US Equities\par}
\vspace{0.3cm}{\large A Comparative Study of Accounting Scores and Learned Return Forecasts\par}
\vspace{0.5cm}{\large Team 34\par}
\vspace{0.15cm}Wan Gehong \quad Wang Tianyi \quad Wu Jiarui \quad Yan Tianao
\par\vspace{0.15cm}12 November 2026
\end{center}
\vspace{0.25cm}
\begin{abstract}
We compare a six-signal accounting rank composite with regularized and tree-based return forecasts in US equities. Two sequential Notebooks prepare raw CRSP and Compustat inputs and study the 2014--2025 panels. The saved execution screens 36 formulaic price--volume candidates and records \FitCount{} new model fits. WQ directions use 2015--2017, WQ screening uses 2018--2019, model cross-validation uses expanding folds in 2020--2022, and portfolio selection uses 2023. Subsequent forecasts are refitted semiannually on the preceding five calendar years. Among twelve score--weighting candidates, the validation rule selects \SelectedName{}. Over 2024--2025 this portfolio earns \SelectedCAGR\% net CAGR, a \SelectedSharpe{} Sharpe ratio, and a \SelectedMDD\% maximum drawdown, against \SPYCAGR\% CAGR for SPY. Returns deduct 10 basis points per dollar bought or sold. The selected portfolio's cost and factor-ablation tests, prediction diagnostics, matched-pool allocation comparisons, and holdings-based nine-style attribution qualify the headline result. Its positive CAGR gap has a paired 95\% interval spanning zero. The risk-model residual is model-dependent and is not certified alpha. This report reads the latest saved results without rerunning the research. The already examined historical sample and latest-download financial vintages limit predictive claims. The principal contribution is a time-ordered, auditable comparison of signals, forecast complexity, allocation, and implementation costs.
\end{abstract}
\vfill\begin{center}\small Keywords: accounting signals; formulaic alphas; purged validation; transaction costs; factor attribution\end{center}
\clearpage
\tableofcontents
\clearpage

\section{Research Question and Status of the Evidence}
Accounting signals are relatively slow-moving, whereas price--volume features react to recent trading. Our economic hypothesis is that the two information sources may be complementary. We test both a transparent accounting ranking and learned combinations, then separate their forecast diagnostics from the performance of actual funded portfolios. The objective is a defensible process rather than a return target.

The direct score retains six financial definitions: four signed-log residual characteristics, free-cash-flow yield, and a seasonal earnings surprise. Ridge and Lasso provide shrinkage-based linear alternatives; shallow XGBoost admits nonlinear interactions. The candidate feature families compare accounting inputs, accounting plus screened WQ inputs, and a larger engineered set under common portfolio rules. Every family contains accounting information: this experiment tests additions to that base, not the incremental value of fundamentals over a pure price--volume model.

\subsection{What changed from the previous report}
The earlier manuscript described a 2023--2025 historical evaluation and a retrospectively emphasized square-root-capitalization portfolio. The present Notebook instead uses 2023 for a declared twelve-candidate portfolio selection and reports 2024--2025 as its subsequent chronological evaluation. It also uses a fixed five-calendar-year forecast-training window, rather than the earlier expanding-from-2015 forecast window. These are substantive research-design changes, not merely new formatting.

\textbf{Crucially, 2024--2025 had already been inspected in earlier research.} Chronological fitting prevents future labels from entering individual forecasts, but it does not make the redesigned experiment a new untouched test. All reported advantages remain retrospective evidence. Neither moving the validation boundary nor repeating a clean run removes specification-search uncertainty.

The evidence is bound to the published 11 November execution record: \DataCodeCells{} preparation cells, \CodeCells{} research cells, \FitCount{} fresh learner fits, and 38 audited portfolio ledgers. Daily financial cross-sectional regressions and direct-score diagnostic calibrations are separate from this count. This 12 November manuscript revision reads the stored outputs; it does not retrain models, replay portfolios, or resample bootstrap intervals. Internal reproducibility does not constitute independent confirmation of profitability.
\clearpage

\section{Data, Coverage, and Independent Preparation}
The preparation Notebook reads raw CRSP daily observations, dated CCM security links, Compustat quarterly fundamentals and annual SIC history, and an archived SPY benchmark response. Monetary accounting is in US dollars. An independent NYSE-session calendar aligns 3,018 sessions during 2014--2025 and 4,366 ever-selected security identifiers; 2014 is feature warm-up.

The monthly universe retains the top 80\% by a trailing 60-session mean of unadjusted close times share volume, after common-equity gates, a \$5 price floor, and 252 observed closing sessions. Daily security and price gates then apply. This yields about two thousand members per date, not 4,366 simultaneously eligible stocks. Three truncated-history checks preserve earlier memberships. Static FIC/PRIUSA descriptors and historical revisions remain limitations. Sina turnover/VWAP coverage is incomplete and is excluded from strategy inputs.
\nbtable{nb_audit}{Saved preparation and source-return audits.}{Raw inputs remain unchanged. The manifest count excludes benchmark and Sina files; snapshot checksums cover the published derived files. Counts are read from saved preparation and strategy outputs.}
\nbfig{nb_universe}{0.95}{Daily reconstructed universe membership. Securities remain on the panel axes outside their eligible dates.}

Preparation quarantines two materially inconsistent output OHLC cells, one inside that day's universe, across prices and trading availability. The material threshold exceeds both \$0.01001 and 0.1\% of closing price. The strategy input audit consequently finds no remaining material violations. Zero volume and economically meaningful negative accounting quantities are retained under their documented rules.
\clearpage

\section{Information Timing, Targets, and Calendar Splits}
Let $P^*_{i,T}$ be the canonical total-return closing chain for stock $i$ on session $T$. Within each known identity segment it is anchored at the first observed raw close and compounded using only past source returns. OHLC ratios and labels retain float64 precision before final feature storage; this avoids premature rounding changes in WQ signs and ranks. The prediction target is
\begin{equation}
y_{i,T}=P^*_{i,T+11}/P^*_{i,T+1}-1.
\end{equation}
Offsets are exchange sessions. Signals formed after $T$ closes execute at the close of $T+1$; the new holding first earns a return on $T+2$. Thus the target contains exactly ten held close-to-close intervals and does not reward same-close execution. It is not a sum of ten simple returns and is not a mandatory holding period.

Targets require finite positive endpoint prices and a consistently known adjustment segment throughout the signal-to-exit interval. Missing exit quotes are not filled. The last eleven signal sessions have no matured target within the snapshot. A training row is admitted only when its exact label-ending date is strictly earlier than its forecast origin. Portfolio returns, unlike prediction diagnostics, retain evaluation dates even when a future target has not yet matured.

\begin{table}[H]\centering\small\caption{Roles of the historical periods.}
\begin{tabularx}{\textwidth}{l X}\toprule Period & Role\\\midrule
2014 & Backward-looking feature warm-up.\\
2015--2017 & Fix WQ directions on matured-target RankIC.\\
2018--2019 & Screen and deduplicate WQ candidates.\\
2020--2022 & Three expanding annual CV folds; two fresh half-year fits per fold.\\
2023 & Compare twelve portfolio specifications using net Sharpe and the cash constraint.\\
2024--2025 & Chronological evaluation of frozen choices; semiannual five-year-window refits.\\\bottomrule
\end{tabularx}\end{table}

Financial fields use supplied availability mappings and forward continuation after their mapped release. Earnings surprise requires its full historical inputs to precede the new release. These timing checks do not reconstruct historical vendor revisions: the financial statements are latest downloaded vintages, not strict PIT observations. No true market-wide VWAP is assumed; labels and execution use adjusted closes, while the implemented WQ subset avoids VWAP formulas rather than renaming close as VWAP.
\clearpage

\section{Six Fundamental Signals}
For dollar-denominated $x$, define $g(x)=\operatorname{sign}(x)\log(1+|x|/10^6)$. Four signals adapt the residual-reconstruction framework in Orient Securities \cite{orient}: retain the residual in a daily eligible cross-sectional OLS regression $g(A)=a+b\,g(B)+\varepsilon$, with at least 30 finite pairs. The million-dollar scaling is part of the US implementation; a signed-log residual is not invariant to arbitrary unit changes.

\begin{table}[H]\centering\small\caption{Accounting definitions and inherited fixed directions.}
\begin{tabularx}{\textwidth}{l X c}\toprule Signal & Definition & Sign\\\midrule
F1 & TTM net-income residual relative to security capitalization & $+1$\\
F2 & Book-equity residual relative to security capitalization & $-1$\\
F3 & Quarterly-sales residual relative to security capitalization & $-1$\\
F4 & Quarterly-net-income residual relative to book equity & $+1$\\
F5 & TTM operating cash flow less capital expenditure, divided by capitalization plus total liabilities less cash & $+1$\\
F6 & Seasonal quarterly earnings surprise, divided by capitalization strictly before release & $+1$\\\bottomrule
\end{tabularx}\end{table}

F5 adapts Northeast Securities' FCF/enterprise-value definition \cite{nesc}, masking nonpositive denominators and SIC2 groups 60--69 at the signal level; financial firms are not removed from the entire study. We do not replicate that report's full quality, liquidity, industry, or portfolio filters. F6 uses China Merchants Securities' simple expectation \cite{cms}: the year-ago quarter's net income plus the average of eight prior year-on-year dollar changes, normalized by pre-release capitalization. The twelve preceding quarters needed for this construction must be available before release.

The six inherited fundamental signs were checked over 2015--2018; the new 2015--2017 direction window applies to WQ, not to a freshly recalibrated accounting family. In particular, negative F2/F3 signs mean the composite does not simply buy more of every accounting residual. Missing components contribute a neutral zero, with the denominator fixed at six.
\nbfig{nb_fundamentals}{0.95}{Signed fundamental RankIC during the 2020--2022 development period. Weak or negative diagnostics are retained rather than removed after inspecting evaluation returns.}
\clearpage

\section{Price--Volume Screening and Feature Families}
We implement 36 supported formulaic candidates from Kakushadze \cite{wq}, not all 101. WQ direction is the sign of mean daily RankIC on 2015--2017 labels matured by that period's end. Screening in 2018--2019 requires signed mean RankIC at least 0.01, coverage at least 50\%, and positive signed RankIC in both screening years. A greedy filter rejects absolute mean daily rank correlation of 0.8 or more, sampled every fifth session. At most ten names are accepted.
\nbtable{nb_wq}{WQ signals retained by the fresh screening run.}{RankIC is an equal-date mean cross-sectional Spearman correlation on finite matured-target pairs; coverage is relative to eligible finite targets. The retained set is not chosen from evaluation returns.}

The accepted factors are \WQList{}. Their short-horizon formulaic vocabulary supplies reversal, trend, price-position, and participation characteristics. Using these definitions with ten-session targets is an adaptation, not a replication of their original published portfolios.

The model compares three feature families. \emph{Fundamental} uses the six financial signals. \emph{Combined} adds the selected WQ values. \emph{Engineered} also adds six 20-session changes in fundamental ranks, 5/20/60-session momentum, 20-session volatility, volume surprise relative to a 20-session mean, and log capitalization. Each value has a missingness indicator. With six retained WQ signals, the largest family has 24 value inputs and 48 total columns, not merely nine model inputs.

All values are converted to eligible same-date percentile ranks centered at 0.5, with missing values set to neutral zero. The missing indicators preserve information about unavailable characteristics. This transformation uses no future fitted scaling moments. All model families share positive-capitalization/common-observability eligibility, so adding a feature does not silently substitute a different evaluation universe.
\clearpage

\section{Learners, Regularization, and Expanding CV}
Ridge offers smooth coefficient shrinkage for noisy correlated characteristics. Lasso permits sparse solutions but may collapse to a constant forecast if penalties overwhelm weak signals. XGBoost allows interactions while restricting depth, shrinkage, and sampling. Complexity is an empirical comparison, not an assumption that trees should win.

Ridge tests four per-sample penalties: $1,0.1,0.01,0.001$; its actual penalty multiplies the coefficient by the training-row count. Lasso tests $0.001,0.0005,0.0001,0.00001,0.000001$, using cyclic coordinate descent, a precomputed Gram matrix, tolerance $10^{-5}$, and at most 5,000 iterations. XGBoost tests depths 2/3/4, learning rates 0.03/0.08, and regularization 1/10, with 200 trees, subsampling 0.8, minimum child weight 100, and seed 5205.

Crossing these grids with three input families produces 12 Ridge, 15 Lasso, and 36 XGBoost configurations: 63 in total. Each has two freshly fitted halves in each of three annual CV folds. There are $63\times6=378$ CV fits, plus 18 validation/evaluation forecast fits, totaling \FitCount{}. We report this full search rather than describing only the final three estimators.

Selection maximizes the equal-fold mean of daily RankIC over 2020--2022. A candidate must converge and have finite correlations on at least 80\% of possible scoring dates in every fold. Ties prefer the earlier, more-shrunk or shallower configuration in the declared ordering. Invalid constant predictions cannot win by an undefined statistic.

\nbtable{nb_hyperparameters}{Frozen feature family and hyperparameters by learner.}{Choices use CV only. Portfolio weighting is selected separately in 2023. CV rows are not independent observations because ten-session labels overlap.}
\clearpage

\section{Forecast Refits and Prediction Integrity}
Within CV, each half-year fit expands from 2015 and uses only targets matured strictly before that half begins. Validation and evaluation change to a fixed preceding five-calendar-year window. At a January or July origin $\tau$, the admissible signal dates satisfy
\[
\tau-5\text{ calendar years}\leq T<\tau,\qquad \operatorname{labelEnd}(T)<\tau.
\]
The comparison between CV and forecast windows is a disclosed modeling convention of this Notebook. This is not the original expanding-from-2015 forecast specification.

\nbtable{nb_windows}{Semiannual forecast training windows, shown for Ridge.}{The other learners use the same date and target-maturity rules. The final label end must strictly precede the forecast start; full training logs are part of the compact evidence.}
\nbfig{nb_cv}{0.88}{Best eligible mean CV RankIC within each feature family and learner. All three scoring folds precede the 2023 portfolio-selection period.}

Every worker records the actual training dates, last label end, sample count, parameter dictionary, convergence status, and fitted feature columns. The fresh verification forbids reused fit checkpoints. Seeds are fixed at 5205 where randomness is used. Ridge/Lasso fits are deterministic conditional on their rows and numerical environment; XGBoost uses a fixed seed and four computational threads.
\clearpage

\section{Validation Selects Both Score and Weighting}
The four score sources are the direct fundamental composite, Ridge, Lasso, and XGBoost. Each is replayed under equal, square-root-capitalization, and capitalization weights. All twelve use the same Top100 target, retention band, schedule, target cap, execution convention, and cost rate. Full-pool allocation diagnostics and SPY are benchmarks, not candidates in the selection grid.

The 2023 rule selects the highest net Sharpe among specifications with finite Sharpe and mean cash at most 5\%. Ties prefer lower turnover, then the stable case name. The winning score and weighting are frozen before the subsequent evaluation fits; a per-model validation choice is also retained for fair presentation. The selected case is \textbf{\SelectedName{}}.
\nbtable{nb_validation}{All twelve 2023 score--weighting candidates.}{An asterisk marks the overall validation selection. High-cash or undefined-Sharpe candidates are ineligible even if another numerical metric appears favorable.}

An ineligible Lasso does not become an investable strategy by breaking ties in identifier order. Constant score dates target cash. Nevertheless, its return-unit forecasts can still have an OOS prediction diagnostic; predictive error and funded-portfolio eligibility are separate questions. No model is dropped solely because its subsequent result is unfavorable.
\clearpage

\section{From Scores to Positions and Trading Costs}
For six signed fundamental characteristics, let $R_T$ be the finite eligible percentile rank and $d_k$ the inherited direction. The direct composite is
\[
z_{ikT}=\begin{cases}R_T(d_k f_{ikT})-0.5,&f_{ikT}\text{ observed},\\0,&\text{otherwise},\end{cases}
\qquad S_{iT}=\frac16\sum_{k=1}^{6}z_{ikT}.
\]
The trading rule uses this ranking, not its later return-unit OLS calibration. At every tenth signal session, incumbents ranked within 150 are retained up to the 100-name target, then remaining slots are filled by score. Seeded SHA256 identifier order resolves boundary ties without dependence on input column order. Next-session quotes never determine replacement selections.

The chosen allocation uses signal-date capitalization, its square root, or equal weights as labeled. Water-filling enforces a 2\% per-name formation target cap. Infeasible residual budget remains cash; drifted or locked holdings can subsequently exceed the formation cap. Portfolios are long-only and unlevered, and cash earns zero.

Let $w^-_{it}$ be beginning-of-period asset weights and $q_t$ absolute traded notional divided by preceding closing wealth. Then
\[
r^{\rm gross}_t=\sum_iw^-_{it}r_{it},\qquad
c_t=\frac{10}{10^4}q_t,\qquad r^{\rm net}_t=r^{\rm gross}_t-c_t.
\]
Both buys and sells are charged. Reported turnover is half the absolute traded notional divided by wealth before trading; it is not substituted directly into the fee formula. Initial entry and attempted terminal liquidation are included. Unfillable buys remain cash, while unsellable holdings remain locked.

Missing held marks are explicitly stale, not economically verified zero returns. Usable terminal source returns take priority; unresolved terminal cases use a recorded last-mark settlement assumption. SPY is an adjusted-close buy-and-hold benchmark over the aligned investment window, with no additional brokerage-fee deduction; strategy and pool ledgers pay the stated costs. Constant-basis-point fees omit market impact, capacity, and time-varying spreads, so these are classroom assumptions rather than executable capacity estimates.
\clearpage

\section{Chronological Evaluation: Return and Risk}
The 2023-selected portfolio is evaluated during 2024--2025 without selecting a replacement from these results. It achieves \SelectedCAGR\% net CAGR, \SelectedVol\% annualized volatility, \SelectedSharpe{} Sharpe, and \SelectedMDD\% maximum drawdown. SPY's aligned CAGR is \SPYCAGR\%. The CAGR gap is \ActiveCAGR{} percentage points; a point estimate alone does not establish statistical or implementable superiority.
\nbtable{nb_evaluation}{Validation-selected representatives and SPY, 2024--2025.}{Each learner's weighting is selected in 2023, not picked from this evaluation table. Ineligible families have no funded representative. Turnover is a daily mean; costs are 10bps per traded dollar.}
\nbfig{nb_wealth}{0.97}{Net wealth of the validation-selected representatives and SPY, with selected-strategy drawdown. Entry and exit fees are included; evaluation holdings are continuous across calendar years.}

The selected Ridge and tree representatives earn \RidgeEvalCAGR\% and \TreeEvalCAGR\% net CAGR. In this experiment, fitting a return predictor does not improve the direct score's funded result. Weak target predictability and trading costs are plausible contributors, but the table does not isolate a causal explanation. The selected score's mean daily turnover is \SelectedTurnover\%.

CAGR compounds net daily returns with a 252-session convention. Annualized volatility uses sample standard deviation times $\sqrt{252}$, and Sharpe uses mean daily return divided by its standard deviation times $\sqrt{252}$, with zero cash interest. Drawdowns include initial wealth in the running peak. These are descriptive conventions, not proof of independent daily returns.
\clearpage

\section{Annual Stability and Comparable Performance Panels}
\nbtable{nb_annual}{Selected portfolio and SPY by period.}{Calendar-year metrics reset the displayed wealth peak to one at the period start; the combined evaluation drawdown uses the continuous two-year path. The 2023 rows are validation evidence, not additional OOS performance.}
\nbfig{nb_periods}{0.98}{Prediction OOS $R^2$, net CAGR, net Sharpe, and maximum drawdown under matched period labels. The scoring-calibration convention differs from the unchanged direct-rank trading rule, as explained next.}

The selected portfolio records \FirstYearCAGR\% CAGR in 2024 against \FirstYearSPY\% for SPY, and \SecondYearCAGR\% in 2025 against \SecondYearSPY\%. Its relative advantage is therefore not stable year by year. The combined positive CAGR gap must be read alongside the larger drawdown. Representatives remain the 2023 choices throughout; we do not choose a different model for each year or chart.
\clearpage

\section{Prediction OOS \texorpdfstring{$R^2$}{R2}: What It Does and Does Not Measure}
We compare squared prediction errors on the same finite eligible stock--day pairs. For each forecast half, the baseline is a constant equal to the mature training targets' historical mean. Targets must mature by the end of the scored phase. For each reported period,
\[
R^2_{\rm OOS}=1-\frac{\sum(y-\widehat y)^2}{\sum(y-\bar y_{\rm train,half})^2}.
\]
The two-year measure pools squared errors rather than averaging annual $R^2$ values. Labels overlap, and millions of stock--days must not be interpreted as millions of independent observations.

The fundamental trading score is not naturally in return units. For diagnostics only, a pooled OLS intercept and slope map it to ten-session returns using eligible data in the preceding five-calendar-year window, with strict target-maturity purging. Both coefficients remain fixed within the forecast half. This calibration never changes the direct portfolio ranking or turns the rank score into a newly fitted trading rule.
\nbtable{nb_r2}{Prediction accuracy against each half's past training-mean forecast.}{Values are decimal $R^2$, not percentages or RankIC. Negative values indicate larger squared error than the historical-mean baseline. The Lasso diagnostic remains visible even when its funded portfolio is ineligible.}

Very small or negative OOS $R^2$ can coexist with a useful cross-sectional ordering because squared error is sensitive to levels and extreme returns, while the portfolio acts on ranks in a selected tail. Conversely, a positive portfolio return does not rescue poor prediction accuracy or prove independent alpha. These are complementary diagnostics with different estimands.
\clearpage

\section{Matched-Universe Allocation and Descriptive Attribution}
\nbtable{nb_weighting}{Identical opportunity sets under alternative weighting rules, 2024--2025.}{The Top100 portfolios have the 2\% formation target cap and retention rule. Full-pool benchmarks hold all eligible positive-capitalization names without that cap; they are not exact reconstructions of SPY.}

Changing weights on the entire pool and changing weights on selected stocks are different experiments. In particular, a loss from square-root weighting on the whole pool cannot be assigned mechanically to the selected portfolio. Selection and weighting interact; a telescoping difference is an accounting identity whose components depend on the chosen comparison path, not a causal decomposition of fundamental skill.

\nbfig{nb_reference_attribution}{0.86}{CAGR-difference attribution for the validation-selected capitalization-weighted strategy. The three path-dependent links sum to its SPY gap; these are descriptive comparisons, not independently tradable or causal premiums.}

The current path is selected capitalization portfolio $\rightarrow$ whole-pool capitalization portfolio $\rightarrow$ whole-pool equal-weight portfolio $\rightarrow$ SPY. Its selection, weighting, and universe links are respectively $+\SelectionLink$, $+\WeightingLink$, and $\UniverseLink$ percentage points, totaling $+\ActiveCAGR$. The first also changes concentration, cap constraints, retention, and costs. The last includes allocation and implementation differences relative to SPY. Neither is pure signal alpha. Reversing or changing the path changes the components.
\clearpage

\section{Robustness: Costs and Factor Ablation}
\nbfig{nb_stress}{0.98}{Saved cost and leave-one-factor-out sensitivity results for the selected six-factor capitalization-weighted strategy. Every scenario retains its 2\% formation cap and other portfolio rules.}

The saved experiment replays the selected ledger at 20 and 30bps, then removes one fundamental contribution at a time and re-averages the remaining five. Sample eligibility, weighting, buffer, and rebalance timing otherwise remain fixed. Costs are computed in the full self-financing ledger, not subtracted linearly from an annual return.

Net CAGR is \SelectedCAGR\% at 10bps, \CostTwentyCAGR\% at 20bps, and \CostThirtyCAGR\% at 30bps, versus \SPYCAGR\% for SPY. The positive historical CAGR gap therefore survives these two higher-cost scenarios. This does not establish statistical robustness or capacity under nonlinear impact.

Dropping F4 raises CAGR to \WithoutFourCAGR\% but lowers Sharpe from \SelectedSharpe{} to \WithoutFourSharpe{}. Thus the results do not show that every factor improves every metric. Removing a signal changes ranks and positions jointly; it is a conditional comparison, not an independent premium estimate. We retain the preselected six-factor rule rather than choosing a subset on these evaluation outcomes. Earlier fusion and financial-delay experiments are outside this latest result set.

Before interpreting net performance as implementable, more realistic execution and data verification would be needed. Constant transaction costs do not capture stock-level spreads, nonlinear impact, order size, or capacity. A low-turnover strategy is cheaper under this model, but low turnover alone does not imply better forecasts or robust outperformance.
\clearpage

\section{Holdings-Based Market, Industry, and Style Attribution}
The risk analysis projects \emph{actual beginning-of-day holdings} of the validation-selected strategy onto a research risk model. Daily weighted cross-sectional regressions explain stock returns using previous-close exposures to Market, SIC industries, and nine styles. Regression weights are proportional to square-root capitalization, and industry returns satisfy a capitalization-weighted zero-sum constraint. This is a US adaptation of the repository's CNE5-style framework, not proprietary MSCI Barra output.

The available styles are Size, Beta, Momentum, Residual Volatility, Nonlinear Size, Book-to-Price, Liquidity, Earnings Yield, and Leverage. Growth is unavailable. Market exposure is invested weight; the standardized Beta style is relative stock beta, not the same quantity. Missing standardized styles enter the portfolio projection at zero and their held weight is reported.
\nbfig{nb_styles}{0.97}{Mean style exposures and Carino-linked cumulative style contributions of the validation-selected portfolio. Exposure units are standardized model units, not portfolio percentages.}
\nbtable{nb_risk_coverage}{Coverage and daily model-fit diagnostics.}{Coverage failures and residual effects remain part of the uncertainty; a numerically reconciled factor projection does not establish that all economic risks have been captured.}

Daily factor contribution is portfolio exposure times estimated factor return. Gross return less systematic contribution is specific/unexplained return. Trading fees remain a separate negative component. Carino linking makes these daily components add to cumulative net return without independently compounding each component as if it were a separate tradable portfolio.
\clearpage

\section{Residual Alpha and Contribution Reconciliation}
\nbtable{nb_components}{Linked cumulative return contributions of the selected portfolio.}{Entries are percentage points of cumulative return over 2024--2025, not annualized returns. Their sum equals cumulative net portfolio profit, including fees.}
The market component contributes \MarketLinked{} cumulative percentage points, the joint style projection \StyleLinked{} points, specific/unexplained return \SpecificLinked{} points, and fees \CostLinked{} points. These conditional projections identify where the fitted model assigns the realized P\&L; they do not identify separately tradable causal premiums.
\nbtable{nb_styles}{Style-level exposures and linked contributions.}{A style's mean exposure and total contribution need not have the same sign because its daily factor return varies. Contributions are conditional projections within this research model.}

Define net residual return as daily net portfolio return minus all projected systematic contributions. Its reported annualized mean is \ResidualAlpha\%, with ten-lag Newey--West $t=\ResidualT$ and a normal-approximation 95\% interval $[\ResidualLow\%,\ResidualHigh\%]$. This estimate differs from the compounded portfolio--SPY gap and from a CAPM intercept.

An interval that excludes zero within this fitted risk model is not proof of a deployable alpha. It does not correct for earlier historical strategy searching, financial revisions, model misspecification, missing industry contributions, stale marks, or terminal-settlement assumptions. Factor returns are estimated contemporaneously for explanatory attribution; they are not additional information made available to the trading strategy.
\clearpage

\section{Paired Uncertainty for the Selected Portfolio}
The selected portfolio and SPY use identical sessions. The saved Notebook reports a paired circular block bootstrap with 2,000 resamples, seed 5205, and 20-session blocks, preserving the paired path and local dependence. Each resample compares independently compounded annualized returns; it does not resample overlapping ten-session labels. We transcribe these existing intervals without new resampling.
\nbtable{nb_intervals}{Saved 20-session paired bootstrap intervals for the current attribution path.}{The Total row is selected strategy minus SPY. Other rows correspond to the Section 13 path. Intervals retain the stored Notebook's displayed precision and are conditional on the research specification; they do not correct for retrospective redesign.}
\nbfig{nb_selected_intervals}{0.90}{The selected-strategy-minus-SPY CAGR gap and its saved paired 95\% interval. Zero lies within the interval.}

With 20-session blocks, the selected portfolio's CAGR-difference interval is $[\SelectedCILow,\SelectedCIHigh]$ percentage points, compared with the point gap of \ActiveCAGR{} points. The interval and the short historical sample should govern interpretation, rather than a standalone positive mean.

The supplementary SPY market regression reports an annualized intercept of 4.22\% with HAC $t=0.786$, distinct from the \ResidualAlpha\% holdings-based risk-model residual. The CAGR gap, regression intercept, and style-model residual are different estimands. A raw-return regression with zero cash interest is a market model, not a risk-free-adjusted CAPM claim. The heuristic $2\,\mathrm{TE}/\sqrt{\mathrm{years}}$ is only an approximate arithmetic-active-return scale, not formal CAGR power analysis.

No confidence interval accounts for every design decision. The same 2024--2025 dates were observed before this redesign. Accordingly, apparent improvements relative to the previous report cannot be interpreted as independent OOS confirmation. A genuinely subsequent sample under a fixed protocol would be required for that claim.
\clearpage

\section{Reproducibility, Interpretation, and Limitations}
Run these entry points in order, each in a fresh kernel:
\begin{quote}\small
\path{notebooks/Equity_Data_Preparation.ipynb}\\
\path{notebooks/Equity_Strategy_Research.ipynb}
\end{quote}
Use the documented Python 3.13 environment and authorized raw sources. Preparation publishes \path{data/processed/latest_snapshot.json} only after success. The strategy Notebook verifies its completion marker, checksums, axes, canonical prices, and label maturity; it stops if no completed input is active. Raw observations remain unchanged.

The stored execution has \DataCodeCells{} data cells and \CodeCells{} research cells in consecutive order, \FitCount{} fresh learner fits, and 38 audited portfolio ledgers. The published audit records 74 passing tests and exact repeated XGBoost predictions on identical input. Independent holdings reconciliation error is at most \texttt{\LedgerError}. A source-return check covers 8,945,343 targets, with maximum absolute error $1.42\times10^{-14}$. These are recorded results, not tests repeated for this manuscript revision.

To rebuild only the report from the current saved outputs, without research execution:
\begin{quote}\small\ttfamily
python scripts/build\_report\_from\_saved\_results.py
\end{quote}
Compile the supplied LaTeX source twice after asset generation. The report inventory binds each saved input to its hash and upstream commit. The wealth-path image is retained from the executed Notebook; aggregate charts are replotted from stored tables. No model is fitted, portfolio replayed, or bootstrap resampled. Restricted data and fitted-model pickles remain outside Git; use trusted team sources for research inputs.

Main limitations are latest financial vintages, release-date proxies, static FIC/PRIUSA descriptors, security-level rather than fully aggregated issuer capitalization, incomplete terminal payouts, constant transaction costs, and historical design reuse. The monthly universe now has an explicit reconstruction and truncation checks, but these do not remove all source-selection uncertainty. The brokerage ideas also postdate part of the evaluation history; this is a retrospective adaptation, not a historically deployed strategy.

The evidence supports a transparent accounting score as the validation-selected candidate in this particular comparison. The richer forecasts do not automatically improve its funded economic result. The practical lesson is to evaluate prediction, portfolio allocation, costs, and uncertainty jointly, while preserving unsuccessful and ineligible alternatives.
\clearpage

\section{Conclusion and Supporting References}
This update integrates data checks, WQ screening, regularized forecasting, portfolio selection, trading costs, and holdings-based risk analysis into one verified execution. The validation rule selects \SelectedName{}, with \SelectedCAGR\% net CAGR and \SelectedSharpe{} Sharpe over the already examined 2024--2025 historical period. The result is not a fresh blind-test discovery, and the conditional risk-model residual is not certified alpha.

We distinguish three questions: whether a model predicts returns in squared-error units, whether its score produces a useful ordering, and whether that ordering survives allocation and costs. Weak OOS $R^2$ does not alone refute a ranking strategy, but neither profitable realized returns nor an explanatory attribution establish robust forecasting ability. A frozen subsequent evaluation and improved data lineage would be needed before a stronger claim.

\begin{thebibliography}{9}
\bibitem{gkx} Gu, S., Kelly, B., and Xiu, D. (2020). Empirical asset pricing via machine learning. \emph{Review of Financial Studies}, 33(5), 2223--2273. \url{https://www.nber.org/papers/w25398}.
\bibitem{wq} Kakushadze, Z. (2016). 101 formulaic alphas. \emph{Wilmott}, 2016(84), 72--80. \url{https://arxiv.org/abs/1601.00991}.
\bibitem{nmv} Novy-Marx, R., and Velikov, M. (2014). A taxonomy of anomalies and their trading costs. NBER Working Paper 20721. \url{https://www.nber.org/papers/w20721}.
\bibitem{orient} Orient Securities Research (21 March 2024). \emph{Reconstructing Fundamental Factors: Factor Stock Selection Series No. 102} [translated title]. Framework p. 8; TTM earnings, sales, book equity and ROE residuals pp. 23, 25, 29, 37. Team-supplied Chinese report; source pages checked against the implemented definitions.
\bibitem{nesc} Northeast Securities Research (22 July 2026, cover date). \emph{Quality and Trend Enhancement of Free-Cash-Flow Strategies: Factor Research No. 5} [translated title]. FCF and enterprise-value definitions p. 5. We adapt the signal, not the full original A-share portfolio.
\bibitem{cms} China Merchants Securities Research (18 May 2026). \emph{Beyond Analyst Expectations: New Explorations in Earnings-Expectation Models} [translated title]. Simple expectations pp. 12--13; pre-announcement-capitalization-normalized earnings surprise p. 27. We adapt the simple-expectation version to the supplied US statements.
\end{thebibliography}

\textbf{Team roster.} Wan Gehong; Wang Tianyi; Wu Jiarui; Yan Tianao.

\textbf{Individual contributions.} Pending confirmation by the four members before submission. Student IDs should be completed by the team if required by the submission template.
\end{document}
